\begin{table}[htbp]\centering
\caption{评测假设与判定标准}\label{tab:hyp}
{\small\begin{threeparttable}
\begin{tabularx}{\linewidth}{lLlL}\toprule
假设 & 对应宣称 & 主要指标 & 判定标准 \\
\midrule
H1 脚步增强有效 & 智能强化关键脚步声；脚步声细节增强 3--9 dB & 脚步增益 G 及其分布 & 人工确认脚步 G 中位数及 95\% CI 下限均 $>$ 3 dB；另报告 G 落在 3--9 dB、$>$ 9 dB、$<$ 3 dB 的比例（描述性，不设通过线） \\
H2 事件选择性（核心） & 实时识别；不同于传统 EQ，只强化必要部分 & 同电平 $\Delta$G、匹配静态 EQ & 同电平枪声 $\Delta$G $\geq$ 3 dB，且匹配静态 EQ 的 |$\Delta$G| $<$ 1 dB \\
H3 非目标声音保持性 & 不干扰环境音或枪声，不改变其他声音特征 & 频谱 MSE、FER & 枪声频谱 MSE 低于实际 EQ，且非脚步 FER（$>$ 3 dB）$\leq$ 10\%（本报告设定的工程容忍阈值，非行业标准） \\
H4 稳定性 & 在目标游戏的复杂战场中表现一致 & 分组重复 H2 & 两款目标游戏、纯净度分档、峰值口径、20\% 标签噪声下，选择性效应方向不变 \\
\bottomrule\end{tabularx}
\end{threeparttable}}
\end{table}
